%!TeX root=../gardentop.tex
\chapter{The Key of the Garden} 
	
\begin{figure}[t!]
\centering
\includegraphics[width=\linewidth]{june}
\end{figure}

 \lettrine[]{T}{wo} days after this, when Mary opened her eyes she sat upright in bed immediately, and called to Martha.

\zz
<Look at the moor! Look at the moor!>

\zz
The rain-storm had ended and the gray mist and clouds had been swept away in the night by the wind. The wind itself had ceased and a brilliant, deep blue sky arched high over the moorland. Never, never had Mary dreamed of a sky so blue. In India skies were hot and blazing; this was of a deep cool blue which almost seemed to sparkle like the waters of some lovely bottomless lake, and here and there, high, high in the arched blueness floated small clouds of snow-white fleece. The far-reaching world of the moor itself looked softly blue instead of gloomy purple-black or awful dreary gray.

<Aye,> said Martha with a cheerful grin. <Th' storm's over for a bit. It does like this at this time o' th' year. It goes off in a night like it was pretendin' it had never been here an' never meant to come again. That's because th' springtime's on its way. It's a long way off yet, but it's comin'.>

<I thought perhaps it always rained or looked dark in England,> Mary said.

<Eh! no!> said Martha, sitting up on her heels among her black lead brushes. <Nowt o' th' soart!>

<What does that mean?> asked Mary seriously. In India the natives spoke different dialects which only a few people understood, so she was not surprised when Martha used words she did not know.

Martha laughed as she had done the first morning.

<There now,> she said. <I've talked broad Yorkshire again like Mrs~Medlock said I mustn't. 'Nowt o' th' soart' means 'nothin'-of-the-sort,'> slowly and carefully, <but it takes so long to say it. Yorkshire's th' sunniest place on earth when it is sunny. I told thee tha'd like th' moor after a bit. Just you wait till you see th' gold-coloured gorse blossoms an' th' blossoms o' th' broom, an' th' heather flowerin', all purple bells, an' hundreds o' butterflies flutterin' an' bees hummin' an' skylarks soarin' up an' singin'. You'll want to get out on it at sunrise an' live out on it all day like Dickon does.>

<Could I ever get there?> asked Mary wistfully, looking through her window at the far-off blue. It was so new and big and wonderful and such a heavenly colour.

<I don't know,> answered Martha. <Tha's never used tha' legs since tha' was born, it seems to me. Tha' couldn't walk five mile. It's five mile to our cottage.>

<I should like to see your cottage.>

Martha stared at her a moment curiously before she took up her polishing brush and began to rub the grate again. She was thinking that the small plain face did not look quite as sour at this moment as it had done the first morning she saw it. It looked just a trifle like little Susan Ann's when she wanted something very much.

<I'll ask my mother about it,> she said. <She's one o' them that nearly always sees a way to do things. It's my day out to-day an' I'm goin' home. Eh! I am glad. Mrs~Medlock thinks a lot o' mother. Perhaps she could talk to her.>

<I like your mother,> said Mary.

<I should think tha' did,> agreed Martha, polishing away.

<I've never seen her,> said Mary.

<No, tha' hasn't,> replied Martha.

She sat up on her heels again and rubbed the end of her nose with the back of her hand as if puzzled for a moment, but she ended quite positively.

<Well, she's that sensible an' hard workin' an' good-natured an' clean that no one could help likin' her whether they'd seen her or not. When I'm goin' home to her on my day out I just jump for joy when I'm crossin' th' moor.>

<I like Dickon,> added Mary. <And I've never seen him.>

<Well,> said Martha stoutly, <I've told thee that th' very birds likes him an' th' rabbits an' wild sheep an' ponies, an' th' foxes themselves. I wonder,> staring at her reflectively, <what Dickon would think of thee?>

<He wouldn't like me,> said Mary in her stiff, cold little way. <No one does.>

Martha looked reflective again.

<How does tha' like thysel'?> she inquired, really quite as if she were curious to know.

Mary hesitated a moment and thought it over.

<Not at all—really,> she answered. <But I never thought of that before.>

Martha grinned a little as if at some homely recollection.

<Mother said that to me once,> she said. <She was at her wash-tub an' I was in a bad temper an' talkin' ill of folk, an' she turns round on me an' says: <Tha' young vixon, tha'! There tha' stands sayin' tha' doesn't like this one an' tha' doesn't like that one. How does tha' like thysel'?> It made me laugh an' it brought me to my senses in a minute.>

She went away in high spirits as soon as she had given Mary her breakfast. She was going to walk five miles across the moor to the cottage, and she was going to help her mother with the washing and do the week's baking and enjoy herself thoroughly.

Mary felt lonelier than ever when she knew she was no longer in the house. She went out into the garden as quickly as possible, and the first thing she did was to run round and round the fountain flower garden ten times. She counted the times carefully and when she had finished she felt in better spirits. The sunshine made the whole place look different. The high, deep, blue sky arched over Misselthwaite as well as over the moor, and she kept lifting her face and looking up into it, trying to imagine what it would be like to lie down on one of the little snow-white clouds and float about. She went into the first kitchen-garden and found Ben Weatherstaff working there with two other gardeners. The change in the weather seemed to have done him good. He spoke to her of his own accord.

<Springtime's comin',> he said. <Cannot tha' smell it?>

Mary sniffed and thought she could.

<I smell something nice and fresh and damp,> she said.

<That's th' good rich earth,> he answered, digging away. <It's in a good humor makin' ready to grow things. It's glad when plantin' time comes. It's dull in th' winter when it's got nowt to do. In th' flower gardens out there things will be stirrin' down below in th' dark. Th' sun's warmin' 'em. You'll see bits o' green spikes stickin' out o' th' black earth after a bit.>

<What will they be?> asked Mary.

<Crocuses an' snowdrops an' daffydowndillys. Has tha' never seen them?>

<No. Everything is hot, and wet, and green after the rains in India,> said Mary. <And I think things grow up in a night.>

<These won't grow up in a night,> said Weatherstaff. <Tha'll have to wait for 'em. They'll poke up a bit higher here, an' push out a spike more there, an' uncurl a leaf this day an' another that. You watch 'em.>

<I am going to,> answered Mary.

Very soon she heard the soft rustling flight of wings again and she knew at once that the robin had come again. He was very pert and lively, and hopped about so close to her feet, and put his head on one side and looked at her so slyly that she asked Ben Weatherstaff a question.

<Do you think he remembers me?> she said.

<Remembers thee!> said Weatherstaff indignantly. <He knows every cabbage stump in th' gardens, let alone th' people. He's never seen a little wench here before, an' he's bent on findin' out all about thee. Tha's no need to try to hide anything from him.>

<Are things stirring down below in the dark in that garden where he lives?> Mary inquired.

<What garden?> grunted Weatherstaff, becoming surly again.

<The one where the old rose-trees are.> She could not help asking, because she wanted so much to know. <Are all the flowers dead, or do some of them come again in the summer? Are there ever any roses?>

<Ask him,> said Ben Weatherstaff, hunching his shoulders toward the robin. <He's the only one as knows. No one else has seen inside it for ten year'.>

Ten years was a long time, Mary thought. She had been born ten years ago.

She walked away, slowly thinking. She had begun to like the garden just as she had begun to like the robin and Dickon and Martha's mother. She was beginning to like Martha, too. That seemed a good many people to like—when you were not used to liking. She thought of the robin as one of the people. She went to her walk outside the long, ivy-covered wall over which she could see the tree-tops; and the second time she walked up and down the most interesting and exciting thing happened to her, and it was all through Ben Weatherstaff's robin.

She heard a chirp and a twitter, and when she looked at the bare flower-bed at her left side there he was hopping about and pretending to peck things out of the earth to persuade her that he had not followed her. But she knew he had followed her and the surprise so filled her with delight that she almost trembled a little.

<You do remember me!> she cried out. <You do! You are prettier than anything else in the world!>

She chirped, and talked, and coaxed and he hopped, and flirted his tail and twittered. It was as if he were talking. His red waistcoat was like satin and he puffed his tiny breast out and was so fine and so grand and so pretty that it was really as if he were showing her how important and like a human person a robin could be. Mistress Mary forgot that she had ever been contrary in her life when he allowed her to draw closer and closer to him, and bend down and talk and try to make something like robin sounds.

Oh! to think that he should actually let her come as near to him as that! He knew nothing in the world would make her put out her hand toward him or startle him in the least tiniest way. He knew it because he was a real person—only nicer than any other person in the world. She was so happy that she scarcely dared to breathe.

The flower-bed was not quite bare. It was bare of flowers because the perennial plants had been cut down for their winter rest, but there were tall shrubs and low ones which grew together at the back of the bed, and as the robin hopped about under them she saw him hop over a small pile of freshly turned up earth. He stopped on it to look for a worm. The earth had been turned up because a dog had been trying to dig up a mole and he had scratched quite a deep hole.

Mary looked at it, not really knowing why the hole was there, and as she looked she saw something almost buried in the newly-turned soil. It was something like a ring of rusty iron or brass and when the robin flew up into a tree nearby she put out her hand and picked the ring up. It was more than a ring, however; it was an old key which looked as if it had been buried a long time.

Mistress Mary stood up and looked at it with an almost frightened face as it hung from her finger.

<Perhaps it has been buried for ten years,> she said in a whisper. <Perhaps it is the key to the garden!>